\section{The Last of the Last Things}

The Church has long taught her children to keep four things in mind: death, judgment, hell, and heaven. The list has the shape of a road, and it is easy to think of the fourth as the last stage of a journey, the town at the end of it. It is not quite that. Death is an event: the debt of nature, paid once by every child of Adam. Judgment is an event: the hour in which what death does not finish is finished, and a life is seen as it is. Hell is a state, the refusal of God made final and suffered for ever, the fire that is not extinguished. But heaven is not one more thing that befalls a creature. It is what the creature is for. The road was built to reach it, and its arrival changes the meaning of every mile behind.

The last of the last things, then, is not a thing. Heaven is a Person, seen. Christ asked it of his Father on the night before he died: ``Father, I will that where I am, they also whom thou hast given me may be with me: that they may see my glory which thou hast given me, because thou hast loved me before the creation of the world.''\endnote{John 17:24, Douay--Rheims (Challoner).} He promised it to the apostles in the same discourse: ``I will come again and will take you to myself: that where I am, you also may be.''\endnote{John 14:3, Douay--Rheims (Challoner).} And the Apostle, writing to Christians who were burying their dead, put the whole hope in a single clause, and drew its use at once: ``and so shall we be always with the Lord. Wherefore, comfort ye one another with these words.''\endnote{1 Thess 4:16--17 (Vulgate numbering), Douay--Rheims (Challoner).}

The year ends where that hope points. The feast of Christ the King sets the Judge on his throne; in Year A of the Postconciliar Missal its Gospel is the scene of the sheep and the goats, where the King says to those on his right hand: ``Come, ye blessed of my Father, possess you the kingdom prepared for you from the foundation of the world.''\endnote{Matt 25:34, Douay--Rheims (Challoner). \emph{Ordo lectionum Missae}, editio typica altera (1981), Domini nostri Iesu Christi universorum Regis, Year A: Matt 25:31--46.} The King who judges is the Face that will be seen; the one who speaks the sentence is the one who prayed that we might be where he is. The last Sundays of the year do not set judgment and heaven against each other. They show that the same Lord is both.

The meditation on judgment and on hell asks for fear, and rightly; fear of the Lord is the beginning of wisdom, and the Church does not let her children forget what may be lost. But the end of fear is love, and the end of the whole meditation on the last things is desire: to want God more than anything that could be lost, and to want him so much that sin loses its taste. The saints who meditated most on death and hell were the same saints who sighed most for heaven. Alphonsus, who wrote the considerations on hell, wrote also: ``Heaven! Heaven!''

\subsection{Following the bier}

The Church does not explain heaven over a Christian's body. When the prayers in the church are finished, says the Roman Ritual, the body is carried to the grave, and while it is carried the clergy sing an antiphon; at the grave the rite goes on, with the blessing of the tomb, the canticle \emph{Benedictus} under the antiphon \latin{Ego sum resurrectio et vita}, and the Lord's Prayer.\endnote{\emph{Rituale Romanum} (Ratisbon: Pustet, 1872), Exsequiarum ordo, pp.~141--143.} The antiphon sung on the way asks for the dead what the Gospel says was given to one poor man. There was a beggar named Lazarus who lay at a rich man's gate, full of sores; ``And it came to pass that the beggar died and was carried by the angels into Abraham's bosom.''\endnote{Luke 16:20, 22, Douay--Rheims (Challoner).} The antiphon asks the angels to lead the soul into paradise, the martyrs to receive it and bring it into the holy city, Jerusalem, and the choir of angels to receive it; and it asks that with Lazarus, once a beggar, the soul may have eternal rest. The holy city it names is the one the Apocalypse saw coming down out of heaven from God.\endnote{Apoc 21:2, Douay--Rheims (Challoner).}

So the meditation on the last things ends where the Church's own steps go: not describing heaven, but praying a soul into it. The prayer is short, because the Church has nothing to add to what she asks. She commits the dead to the angels and the martyrs, and asks for them the holy city and eternal rest. And she lets the living, following the bier, hear what they too are made for.
